% ============================================================
%  Regulación del Alquiler y Dinámica del Mercado Residencial
%  en Barcelona (2000–2025): Análisis Multimétodo con Panel,
%  Clasificación Supervisada y Series Temporales
%
%  Compilar en Overleaf (pdfLaTeX)
% ============================================================
\documentclass[12pt,a4paper]{article}

% ── Encoding & idioma ──────────────────────────────────────
\usepackage[utf8]{inputenc}
\usepackage[T1]{fontenc}
\usepackage[spanish,english]{babel}

% ── Márgenes ───────────────────────────────────────────────
\usepackage[top=2.5cm, bottom=2.5cm, left=2.8cm, right=2.8cm]{geometry}

% ── Tipografía ─────────────────────────────────────────────
\usepackage{microtype}
\usepackage{lmodern}

% ── Matemáticas ────────────────────────────────────────────
\usepackage{amsmath, amssymb}
\usepackage{bm}

% ── Tablas ─────────────────────────────────────────────────
\usepackage{booktabs}
\usepackage{threeparttable}
\usepackage{multirow}
\usepackage{array}
\usepackage{lscape}

% ── Figuras ────────────────────────────────────────────────
\usepackage{graphicx}
\usepackage{subcaption}
\usepackage{float}
\usepackage[dvipsnames]{xcolor}

% ── Hipervínculos y referencias ────────────────────────────
\usepackage[colorlinks=true, linkcolor=NavyBlue,
            citecolor=NavyBlue, urlcolor=NavyBlue]{hyperref}
\usepackage{natbib}

% ── Espaciado ──────────────────────────────────────────────
\usepackage{setspace}
\onehalfspacing

% ── Código inline ──────────────────────────────────────────
\usepackage{listings}
\lstset{basicstyle=\small\ttfamily, breaklines=true}

% ── Cabeceras y pies ───────────────────────────────────────
\usepackage{fancyhdr}
\pagestyle{fancy}
\fancyhf{}
\rhead{\small Regulación del alquiler en Barcelona · 2000--2025}
\lhead{\small La Salle -- Universitat Ramon Llull}
\cfoot{\thepage}
\renewcommand{\headrulewidth}{0.4pt}

% ── Otros ──────────────────────────────────────────────────
\usepackage{url}
\usepackage{enumitem}
\usepackage{caption}
\captionsetup{font=small, labelfont=bf}

% ============================================================
\begin{document}
\selectlanguage{spanish}

% ── PORTADA ─────────────────────────────────────────────────
\begin{titlepage}
  \centering
  \vspace*{1cm}

  {\Large \textbf{La Salle -- Universitat Ramon Llull}\\[0.3em]
  \large Máster en Inteligencia Artificial / Data Science}

  \vspace{1.5cm}

  {\LARGE \bfseries
   Regulación del Alquiler y Dinámica del Mercado\\[0.4em]
   Residencial en Barcelona (2000--2025):\\[0.5em]
   Un Análisis Multimétodo con Panel de Datos,\\[0.4em]
   Clasificación Supervisada y Series Temporales}

  \vspace{0.8cm}
  {\large \textit{Rental Regulation and Residential Market Dynamics in\\
  Barcelona (2000--2025): A Multi-Method Analysis with\\
  Panel Data, Supervised Classification and Time Series}}

  \vspace{1.5cm}

  {\large \textbf{Autores:}\\[0.4em]
  Katia Soto Huerta\\[0.2em]
  \texttt{kthsoto@gmail.com}\\[0.4em]
  La Salle -- Universitat Ramon Llull, Barcelona}

  \vspace{1.2cm}
  {\large Abril de 2026}

  \vfill
  \hrule
  \vspace{0.3cm}
  \small{Proyecto 2.1 · Mini-TFM · Área: Econometría aplicada y Machine Learning}
\end{titlepage}

% ── RESUMEN ─────────────────────────────────────────────────
\begin{abstract}
\noindent
\textbf{Resumen.}
El mercado de alquiler de Barcelona ha atravesado tres ciclos regulatorios entre 2020 y 2025
(Ley 11/2020, derogación del Tribunal Constitucional y Ley 12/2023), ofreciendo una ventana
natural de análisis cuasi-experimental. Este trabajo combina cuatro métodos cuantitativos sobre
datos oficiales trimestrales de 10 distritos y 66 barrios (2000Q1--2025Q3): \textit{(i)} regresión
de panel con efectos fijos (2.838 obs.); \textit{(ii)} clasificación supervisada de episodios de
tensión contractual por barrio; \textit{(iii)} series temporales SARIMAX$(1,1,1)(1,0,1)_4$ con
análisis de intervención ITS-3 a nivel distrito; y \textit{(iv)} los mismos modelos sobre la serie
agregada ciudad en tres métricas: precio nominal, precio real (base 2025) y volumen de contratos.
El ITS-3 muestra que la Ley 12/2023 produce una caída significativa en el número de contratos
($\hat{\beta}_{D_{ley12}} = -1.588$, $p = 0.007$) pero un incremento positivo en el precio real
(+0.495 \euro/m\textsuperscript{2}, $p = 0.034$), coherente con una contracción de la oferta al
segmento barato del mercado. El SARIMAX alcanza un MAPE del 3.3\,\% en el precio nominal agregado.
\end{abstract}

\vspace{0.4cm}
\noindent
\textbf{Palabras clave:} regulación del alquiler, análisis de series temporales interrumpidas,
datos de panel con efectos fijos, SARIMAX, tensión contractual, zona tensionada, Barcelona,
efecto composición, econometría aplicada.

\vspace{0.4cm}

% Abstract en inglés
\begin{otherlanguage}{english}
\begin{abstract}
\noindent
\textbf{Abstract.}
Barcelona's rental market has undergone three regulatory cycles between 2020 and 2025
(Law 11/2020, Constitutional Court ruling, Law 12/2023), providing a natural quasi-experimental
window of analysis. This paper combines four quantitative methods on official quarterly data for
10 districts and 66 neighbourhoods (2000Q1--2025Q3): \textit{(i)} fixed-effects panel regression
(2,838 obs.); \textit{(ii)} supervised classification of contractual tension episodes at
neighbourhood level; \textit{(iii)} SARIMAX$(1,1,1)(1,0,1)_4$ time-series models with
multi-break ITS analysis at district level; and \textit{(iv)} the same models on the city-level
aggregate series for three metrics: nominal price, real price (2025 base), and contract volume.
The ITS-3 shows that Law 12/2023 produces a significant drop in contract volume
($\hat{\beta}_{D_{ley12}} = -1{,}588$, $p = 0.007$) but a positive shift in real price
(+0.495 \euro/m\textsuperscript{2}, $p = 0.034$), consistent with a supply contraction in the
lower-price market segment. The SARIMAX achieves a 3.3\,\% MAPE on the city-level nominal price
series.\\[0.3em]
\textbf{Keywords:} rent control, interrupted time series, panel data, fixed effects, SARIMAX,
contractual tension, tensioned market, Barcelona, composition effect, applied econometrics.
\end{abstract}
\end{otherlanguage}

\newpage
\tableofcontents
\newpage

% ============================================================
\section{Introducción}
\label{sec:intro}

La vivienda es uno de los principales focos de tensión socioeconómica en las grandes
ciudades europeas del siglo XXI. El acceso al alquiler asequible se ha convertido en
un problema estructural que trasciende el ciclo económico y plantea retos tanto de
política pública como de análisis cuantitativo. Barcelona representa un caso
especialmente relevante: en el período 2020--2025 ha sido escenario de tres ciclos
regulatorios sucesivos en un lapso de apenas cinco años, lo que proporciona una
variación temporal exógena inusualmente rica para el análisis estadístico.

\subsection*{Motivación e importancia de la aportación}

La literatura académica sobre regulación del alquiler ha crecido sustancialmente en la
última década \citep{diamond2019rent, sims2007rent}, pero la mayor parte de los trabajos
empíricos se han centrado en mercados anglosajones, con datos seccionales o en mercados
donde solo existe un único evento regulatorio. El caso de Barcelona ofrece dos ventajas
metodológicas adicionales: \textit{(i)} un horizonte temporal de 25 años con datos
trimestrales desagregados por distrito y barrio, y \textit{(ii)} tres rupturas
institucionales de naturaleza distinta (regulación, derogación, re-regulación), lo que
permite comparar la magnitud y dirección del efecto según el tipo de intervención.

Frente a los estudios anteriores para España \citep{garcia2018housing, modenes2021alquiler},
este trabajo aporta una perspectiva multimétodo que triangula econometría de panel,
clasificación supervisada y series temporales sobre \emph{tres métricas} del mercado
simultáneamente: precio nominal, precio real deflactado y volumen contractual. Esta
triangulación permite distinguir si un cambio observado es puramente monetario, un
reajuste estructural del precio real o una contracción de la oferta.

\subsection*{Marco de trabajo}

El análisis se estructura en cuatro módulos complementarios que responden a preguntas
distintas sobre el mismo fenómeno:

\begin{enumerate}[leftmargin=*, label=\textbf{M\arabic*.}]
  \item \textbf{Regresión de panel (FE):} ¿Tiene la regulación un efecto medio
    significativo sobre el crecimiento de contratos, una vez controlados precio,
    superficie y tipo de interés?
  \item \textbf{Clasificación supervisada:} ¿Es predecible qué barrios entrarán
    en episodios de tensión contractual bajo el nuevo régimen?
  \item \textbf{Series temporales (nivel distrito):} ¿Cuánto mejora el SARIMAX
    sobre los baselines en cada combinación (distrito, métrica)?
    ¿Qué magnitud y significancia tienen los cambios estructurales en cada
    ruptura regulatoria, desagregados por distrito y métrica?
  \item \textbf{Series temporales (nivel ciudad):} Validación de los resultados
    anteriores sobre el agregado city-level con proyecciones condicionales.
\end{enumerate}

\subsection*{Estructura del artículo}

La Sección~\ref{sec:relacionados} revisa la literatura relevante.
La Sección~\ref{sec:antecedentes} describe el contexto regulatorio.
La Sección~\ref{sec:tecnologias} presenta las tecnologías empleadas.
La Sección~\ref{sec:propuesta} desarrolla la metodología de cada módulo.
La Sección~\ref{sec:resultados} reporta los resultados experimentales.
La Sección~\ref{sec:discusion} abre la discusión cualitativa.
La Sección~\ref{sec:conclusiones} cierra con conclusiones y líneas abiertas.

% ============================================================
\section{Trabajo Relacionado y Antecedentes}
\label{sec:relacionados}

\subsection{Efectos de la regulación del alquiler}

El debate académico sobre el control de rentas tiene una larga historia. Los modelos
de primera generación \citep{arnott1995rent} predicen que cualquier tope de precio
por debajo del equilibrio reduce la oferta y deteriora la calidad del parque. La
evidencia empírica más sólida proviene de \citet{diamond2019rent}, que aprovechan
una ampliación del control de rentas en San Francisco como experimento natural y
documentan una caída del 15\,\% en la oferta de viviendas sometidas a regulación en
el largo plazo, con una recomposición hacia usos alternativos (demolición, conversión
a propiedad). \citet{sims2007rent} llegan a conclusiones similares analizando el fin
del control de rentas en Massachusetts.

Desde una perspectiva europea, \citet{jenkins2009review} revisa la evidencia de
múltiples países y concluye que la regulación de segunda generación ---que permite
actualizaciones moderadas vinculadas al IPC--- produce efectos menos distorsionadores
que los topes estrictos, aunque mantiene el riesgo de contracción de la oferta nueva.
El caso de Berlín con la \textit{Mietendeckel} (2020--2021) ofrece un paralelo
directo: estudios de diferencias en diferencias encuentran caídas de precio en el
segmento regulado pero reducción paralela de la oferta y aumento de precio en el
segmento no regulado \citep{kholodilin2021rent}.

\subsection{Mercado de alquiler en España y Cataluña}

La literatura específica para España es más escasa. \citet{garcia2018housing}
documentan la deterioración de la asequibilidad en Barcelona entre 2014 y 2018,
relacionándola con la recuperación económica y el auge del turismo. \citet{modenes2021alquiler}
analiza el cambio estructural del régimen de tenencia en España durante la crisis
financiera, con el alquiler creciendo del 14\,\% al 23\,\% de los hogares entre 2008
y 2020. Para la Ley 11/2020, los primeros análisis descriptivos de \citet{bosch2021ley11}
muestran una contracción de la oferta en los meses inmediatos a su entrada en vigor,
pero el horizonte temporal disponible era muy limitado.

\subsection{Métodos cuantitativos en análisis de políticas de vivienda}

El diseño de series temporales interrumpidas (ITS) de \citet{wagner2002its} es el
estándar metodológico para evaluar intervenciones de política cuando no existe un
grupo de control equivalente. La especificación ITS con múltiples rupturas permite
descomponer el efecto total en cambios de nivel (\textit{level shift}) y cambios de
pendiente (\textit{slope change}) para cada intervención, controlando por tendencia
preexistente y variables exógenas. Los errores HAC de Newey-West \citep{newey1987hac}
son imprescindibles en este contexto para corregir la autocorrelación serial de los
residuos OLS.

Para la componente de pronóstico, seguimos la especificación SARIMAX de \citet{box2015timeseries}
y el marco de evaluación de \citet{hyndman2021forecasting}, que recomienda el uso de
baselines estacionales como referencia mínima antes de aceptar modelos más complejos.
En regresión de panel, el estimador \textit{within} con errores clustered es el
estándar de \citet{baltagi2021econometric}; la justificación de la identificación
cuasi-experimental sigue a \citet{angrist2009mostly}.

El uso de Random Forest y Gradient Boosting para clasificación de riesgo en mercados
inmobiliarios es relativamente reciente pero creciente \citep{genesove2012behavioral}.
La métrica ROC-AUC con validación cruzada temporal (\textit{TimeSeriesSplit}) es el
protocolo recomendado cuando la variable objetivo tiene dependencia temporal, para
evitar data leakage \citep{bergmeir2018note}.

% ============================================================
\section{Antecedentes: Marco Regulatorio de Barcelona}
\label{sec:antecedentes}

El análisis distingue cinco períodos regulatorios bien delimitados (Tabla~\ref{tab:periodos}).

\begin{table}[H]
\caption{Períodos regulatorios estudiados}
\label{tab:periodos}
\centering
\small
\begin{threeparttable}
\begin{tabular}{clll}
\toprule
Código & Período & Fechas & Descripción \\
\midrule
P1 & Pre-COVID / pre-regulación & 2000Q1--2019Q4 & Mercado libre \\
P2 & COVID shock                & 2020Q1--2020Q3 & Confinamiento; caída contratos $-22$\,\% \\
P3 & Ley 11/2020 activa         & 2020Q4--2022Q1 & Topes de renta (Catalunya) \\
P4 & Desregulación TC           & 2022Q2--2023Q2 & TC derogó Ley 11/2020; mercado libre \\
P5 & Ley 12/2023 / Z.\ Tensionada& 2023Q3--hoy   & Zona tensionada estatal; índice referencia \\
\bottomrule
\end{tabular}
\end{threeparttable}
\end{table}

La \textbf{Ley 11/2020} del Parlament de Catalunya fue la primera regulación de rentas
de alquiler en España desde la transición democrática. Introducía un índice de
referencia de precio de alquiler y un tope a la renta en zonas declaradas tensionadas.

En \textbf{marzo de 2022}, el Tribunal Constitucional declaró inconstitucional el
artículo 6 (topes de renta) por invasión de competencias estatales, dejando Barcelona
temporalmente sin regulación de precios. Este período de vacío regulatorio (P4) es
especialmente relevante porque ofrece una ventana de identificación natural.

La \textbf{Ley 12/2023 de Vivienda} (estatal) volvió a declarar Barcelona zona de
mercado residencial tensionado en julio de 2023, estableciendo un nuevo índice de
referencia y recuperando el espíritu de la Ley 11/2020 bajo cobertura constitucional.

% ============================================================
\section{Tecnologías Utilizadas}
\label{sec:tecnologias}

El análisis se implementa íntegramente en Python\,3.11 sobre entorno Jupyter Notebook.
La Tabla~\ref{tab:tech} resume las tecnologías principales.

\begin{table}[H]
\caption{Stack tecnológico del análisis}
\label{tab:tech}
\centering
\small
\begin{tabular}{lll}
\toprule
Componente & Librería / Fuente & Uso principal \\
\midrule
Manipulación datos  & \texttt{pandas}, \texttt{numpy}             & ETL, índices temporales \\
Panel econométrico  & \texttt{linearmodels} 6.x                  & PanelOLS, within estimator \\
Series temporales   & \texttt{statsmodels} 0.14.x                & SARIMAX, ETS, ADF, KPSS, ITS \\
Clasificación       & \texttt{scikit-learn} 1.4.x                & RF, GB, LR, ROC-AUC \\
Visualización       & \texttt{matplotlib}, \texttt{seaborn}       & Gráficas y diagnósticos \\
Datos fuente        & Portal Dades Obertes Ajuntament BCN         & Contratos lloguer 2000--2025 \\
Deflactor           & IPC INE (base 2025)                         & Precio real \\
Variable macro      & Euríbor 12m (Banco de España)               & Exógena modelos \\
\bottomrule
\end{tabular}
\end{table}

El dataset principal (\textit{Contractes de lloguer per districtes i barris}) se
descarga como Parquet desde el portal de datos abiertos del Ajuntament de Barcelona.
Los macroeconómicos (IPC, Euríbor) se añaden por \textit{merge} temporal. El entorno
completo es reproducible mediante \texttt{requirements.txt}.

% ============================================================
\section{Desarrollo de la Propuesta}
\label{sec:propuesta}

\subsection{El problema}

El mercado de alquiler residencial barcelonés presenta una complejidad que dificulta
el análisis con un único método: el precio y el volumen responden a la regulación con
intensidad y dirección distintas, la heterogeneidad geográfica es sustancial (Nou
Barris vs.\ Eixample), y la serie temporal cubre veinticinco años con tres rupturas
institucionales y un shock pandémico. Un análisis que solo mire precios puede
concluir que la regulación ``no funcionó'' cuando en realidad la señal se ha
desplazado al mercado de contratos de temporada (no observables en los datos). Un
análisis que solo mire volumen puede exagerar el efecto si parte de la caída se debe
al COVID o a la subida del Euríbor.

La estrategia multimétodo que proponemos aborda esta complejidad desde cuatro ángulos
complementarios, cada uno con su diseño de identificación propio.

\subsection{El dataset}

Los datos provienen del Portal de Dades Obertes del Ajuntament de Barcelona,
\textit{dataset}: Habitatge --- Contractes de lloguer per districtes i barris (actualización
trimestral). La cobertura es 2000Q1--2025Q3 (103 trimestres), con información para
10 distritos y 73 barrios, de los cuales se excluyen 7 de mercado delgado
($< 300$ contratos medianos por trimestre) para garantizar robustez estadística.

Las variables analíticas principales son:
\begin{itemize}[leftmargin=*, itemsep=2pt]
  \item \texttt{avg\_rent\_m2}: precio medio del alquiler en \euro/m\textsuperscript{2}
    (nominal), construido como media ponderada por \texttt{num\_contracts} de los
    submercados observados.
  \item \texttt{avg\_rent\_m2\_real\_2025base}: precio deflactado con el IPC,
    base 2025, disponible desde 2002Q1.
  \item \texttt{num\_contracts}: número de contratos registrados por trimestre.
  \item \texttt{avg\_surface}: superficie media contratada (m\textsuperscript{2}).
  \item \texttt{euribor\_12m\_q}: Euríbor a 12 meses, media trimestral.
  \item Dummies regulatorias: \texttt{post\_regulation}, \texttt{covid\_dummy},
    $D_{ley11}$, $D_{tc}$, $D_{ley12}$ según la periodización de la
    Tabla~\ref{tab:periodos}.
\end{itemize}

El pipeline de limpieza incluye: recodificación de ceros anómalos como NaN, construcción
del índice temporal completo \texttt{QS-OCT} con interpolación temporal (\texttt{method='time'})
para huecos internos, y agregación ciudad mediante media ponderada por número de contratos.
El panel analítico final tiene $N=66$ barrios, $T=103$ trimestres y 2.838 observaciones
válidas para los modelos de crecimiento.

\subsection{La aplicación: cuatro módulos metodológicos}

\subsubsection{Módulo 1 --- Regresión de Panel con Efectos Fijos}

La variable dependiente es la tasa de crecimiento interanual de contratos
(\texttt{contract\_growth\_yoy}). Estimamos el modelo:
\begin{equation}
  y_{it} = \alpha_i + \beta_1 \,\text{post\_regulation}_{it}
         + \beta_2 \,\text{covid}_{it}
         + \gamma'\mathbf{x}_{it} + \varepsilon_{it}
  \label{eq:panel}
\end{equation}
donde $i$ indexa el barrio, $t$ el trimestre, $\alpha_i$ son efectos fijos barrio
(estimador \textit{within}), y $\mathbf{x}_{it}$ incluye \texttt{avg\_rent\_m2},
\texttt{avg\_surface}, \texttt{euribor\_12m\_q} y dummies estacionales $q_2, q_3, q_4$.
Los errores estándar se estiman con clustering por barrio para ser robustos a
heterocedasticidad y autocorrelación intra-cluster \citep{baltagi2021econometric}.

El estimador within elimina toda heterogeneidad barrio-específica fija, de modo que
la identificación proviene de la variación \textit{within}-barrio a lo largo del tiempo.
La especificación completa (M3) añade el precio y el Euríbor para controlar los
determinantes de mercado más relevantes.

Para el análisis de heterogeneidad temporal, el modelo M4 reemplaza
\texttt{post\_regulation} por tres indicadores de régimen:
\begin{equation}
  y_{it} = \alpha_i + \delta_1 D_{ley11,t} + \delta_2 D_{tc,t} + \delta_3 D_{ley12,t}
         + \gamma'\mathbf{x}_{it} + \varepsilon_{it}
  \label{eq:panel_m4}
\end{equation}

\subsubsection{Módulo 2 --- Clasificación de Tensión Contractual}

Definimos un episodio de \textit{tensión contractual} para el barrio $i$ en el
trimestre $t$ como:
\begin{equation}
  \text{market\_tension}_{it} =
  \mathbf{1}\!\left[
    \text{cg}_{it} < \hat{q}_{25,i}^{(\leq 2019)}
    \;\wedge\;
    \text{cg}_{it} < -0.05
  \right]
  \label{eq:tension}
\end{equation}
donde $\hat{q}_{25,i}^{(\leq 2019)}$ es el percentil 25 de crecimiento del barrio $i$
calculado sobre el período pre-regulación. El umbral del $-5$\,\% evita que variaciones
menores activen la señal. Entrenamos tres clasificadores (Regresión Logística L2,
Random Forest, Gradient Boosting) con un split temporal estricto: \textit{train}
$\leq 2021$, \textit{test} $\geq 2022$. La evaluación principal es el ROC-AUC,
complementado con validación cruzada temporal de 5 folds (\texttt{TimeSeriesSplit}).

\subsubsection{Módulo 3 --- SARIMAX e ITS a Nivel Distrito}

Para cada combinación (distrito, métrica) --- $10 \times 3 = 30$ series --- ajustamos
la especificación:
\begin{equation}
  \phi(B)\,\Phi(B^4)\,(1-B)\,y_{d,t}
  = c_d + \theta(B)\,\Theta(B^4)\,\varepsilon_{d,t}
    + \gamma_d\,\text{euribor}_t
  \label{eq:sarimax}
\end{equation}
con orden $(1,1,1)(1,0,1)_4$, frecuencia \texttt{QS-OCT} y estimación por MLE con
\texttt{cov\_type='oim'}. El orden $d=1$ queda justificado por la convergencia ADF+KPSS
en 29 de 30 combinaciones (Tabla~\ref{tab:stat_orden}). La exógena única
\texttt{euribor\_12m\_q} se elige por estabilidad numérica: las dummies regulatorias
son idénticas en todos los distritos y son cero en todo el período de entrenamiento
(hasta 2019Q4), produciendo singularidad si se incluyen.

Los cortes son: \textit{train} $< 2020$Q1, \textit{test baselines} $\geq 2022$Q1 y
\textit{test SARIMAX} $\geq 2020$Q1 (para cubrir el período COVID que el modelo
explícitamente modela con el Euríbor).

El ITS-3 multi-ruptura se especifica como:
\begin{equation}
  y_t^{(m)} = \beta_0 + \beta_1 t
  + \sum_{k \in \{11, tc, 12\}}
      \!\bigl[\beta_{2k} D_{k,t} + \beta_{3k}\, t_{post_k,t}\bigr]
  + \beta_4\,\text{covid}_t + \beta_5\,\text{euribor}_t + \varepsilon_t
  \label{eq:its3}
\end{equation}
donde $D_{k,t} = \mathbf{1}[t \geq t_k^*]$ es el escalón (cambio de nivel inmediato) y
$t_{post_k,t} = \max(0,\, t - t_k^*)$ es la rampa post-ruptura (cambio de pendiente).
Los errores se estiman con HAC Newey-West ($L = 4$ lags) \citep{newey1987hac}.

\subsubsection{Módulo 4 --- Series Temporales Ciudad Agregada}

Replicamos los modelos del Módulo 3 sobre la serie ciudad construida como media
ponderada por \texttt{num\_contracts} de los 10 distritos. Esta capa proporciona una
validación de coherencia: si los signos y magnitudes de los coeficientes ITS-3 ciudad
son consistentes con la mediana de los coeficientes distrito, la narrativa global es
robusta; si divergen, existe heterogeneidad geográfica oculta.

\subsection{Diseño experimental}

\paragraph{Panel.} Split temporal no aleatorio: la identificación proviene de la
variación within-barrio. Robustez con precio real (\textbf{M\_ROB}).

\paragraph{Clasificación.} Split estrictamente temporal (train $\leq 2021$, test
$\geq 2022$) para simular la predicción en tiempo real. Validación cruzada con
\texttt{TimeSeriesSplit} de 5 folds. Sensibilidad del umbral $p_{25}$ variando entre
$p_{20}$ y $p_{30}$ (\textbf{robustez de la variable objetivo}).

\paragraph{Series temporales.} Los baselines naïve estacional y ETS establecen el
suelo de rendimiento. MAPE y RMSE sobre test compartido (2022Q1+). Diagnóstico de
residuos con Ljung-Box y Durbin-Watson. Sensibilidad del ITS-3 a la fecha exacta
del punto de corte (ventanas de $\pm 2$ trimestres).

% ============================================================
\section{Experimentación y Resultados}
\label{sec:resultados}

\subsection{EDA y cambio estructural}

La Tabla~\ref{tab:prepost} muestra las diferencias estadísticas entre períodos
pre y post-regulación (2020Q4 en adelante) para las cuatro variables principales.

\begin{table}[H]
\caption{Tests pre/post-regulación (Welch y Mann-Whitney)}
\label{tab:prepost}
\centering
\small
\begin{threeparttable}
\begin{tabular}{lrrrrl}
\toprule
Variable & $\Delta$\,Media\,(\%) & $p$-Welch & $p$-MWU & Sig. \\
\midrule
\texttt{num\_contracts}             & $+22.9$ & $0.0056$ & $0.0085$ & ** \\
\texttt{avg\_rent\_m2}              & $+41.7$ & $<0.001$ & $<0.001$ & *** \\
\texttt{avg\_rent\_m2\_real\_2025}  & $+7.2$  & $<0.001$ & $0.0066$ & ** \\
\texttt{avg\_surface}               & $+3.0$  & $<0.001$ & $<0.001$ & *** \\
\bottomrule
\end{tabular}
\begin{tablenotes}
\small\item Nota: $n_{pre}=83$, $n_{post}=20$ observaciones trimestrales ciudad.
$+41.7$\,\% nominal frente a $+7.2$\,\% real refleja el componente inflacionario.
$^{***}p<0.01$, $^{**}p<0.05$.
\end{tablenotes}
\end{threeparttable}
\end{table}

El test de Chow simplificado confirma ruptura estructural en 2020Q4 para el precio
nominal ($F = 5.16$, $p = 0.007$) y el número de contratos ($F = 101.2$, $p < 0.001$),
pero no para el precio real ($F = 0.28$, $p = 0.755$), lo que indica que la ruptura
observada en nominal es en gran medida inflacionaria.

La heterogeneidad entre distritos es sustancial: Sant Martí y Eixample lideran la
variación de renta real post-regulación ($+10.8$\,\% y $+10.3$\,\% respectivamente),
mientras Nou Barris prácticamente no varía ($+0.7$\,\%). Esta dispersión justifica
la elección del estimador within y el análisis desagregado a nivel distrito.

\subsection{Módulo 1: Regresión de Panel}

La Tabla~\ref{tab:panel} presenta la comparativa de modelos M1--M3.

\begin{table}[H]
\caption{Comparativa de especificaciones Panel FE (variable dep.: \texttt{contract\_growth\_yoy})}
\label{tab:panel}
\centering
\small
\begin{threeparttable}
\begin{tabular}{lccc}
\toprule
 & M1 Baseline & M2 Precio+Sup & M3 Completo \\
\midrule
\texttt{post\_regulation}
  & $-0.015$ & $0.136^{***}$ & $0.231^{***}$ \\
  & $(-1.89)$ & $(9.96)$ & $(15.71)$ \\
\texttt{covid\_dummy}
  & $0.204^{***}$ & $0.159^{***}$ & $0.040^{**}$ \\
  & $(13.36)$ & $(10.03)$ & $(2.29)$ \\
\texttt{avg\_rent\_m2}
  & --- & $-0.066^{***}$ & $-0.038^{***}$ \\
  & & $(-10.78)$ & $(-4.86)$ \\
\texttt{avg\_surface}
  & --- & $-0.005$ & $-0.003$ \\
\texttt{euribor\_12m\_q}
  & --- & --- & $-0.076^{***}$ \\
  & & & $(-7.96)$ \\
Efectos estacionales
  & No & No & Sí \\
\midrule
$R^2$ Within & $0.029$ & $0.089$ & $0.128$ \\
Obs.\ & $2{,}838$ & $2{,}838$ & $2{,}838$ \\
\bottomrule
\end{tabular}
\begin{tablenotes}
\small\item \textit{t}-estadísticos entre paréntesis. SE clustered por barrio.
$^{***}p<0.01$, $^{**}p<0.05$, $^{*}p<0.10$.
Todos los modelos incluyen efectos fijos por barrio (estimador within).
\end{tablenotes}
\end{threeparttable}
\end{table}

El cambio de signo de \texttt{post\_regulation} entre M1 ($-0.015$, ns) y M3
($+0.231$, $p<0.001$) se explica por la mediación del precio: al controlar por
\texttt{avg\_rent\_m2}, emerge el efecto directo de la regulación sobre la rotación
contractual del segmento que \textit{permanece} en el mercado observable. No es
paradójico: los propietarios que no retiran su vivienda bajo regulación exhiben mayor
rotación de arrendatarios. El efecto nivel (precio) es el que captura la presión real.

El coeficiente del Euríbor en M3 ($-0.076$, $p<0.001$) confirma el canal
financiero: subidas del tipo de interés se asocian a caídas en el ritmo contractual.

\paragraph{Modelo M4 --- tres regímenes.} Los coeficientes son: $\hat{\delta}_1 =
+0.317$ (Ley 11/2020, $p<0.001$), $\hat{\delta}_2 = +0.001$ (derogación TC, ns) y
$\hat{\delta}_3 = +0.028$ (Ley 12/2023, ns). La significancia solo del primer régimen
en el estimador within refleja que el shock de 2020Q4 (rebote post-COVID + nueva
regulación) domina la señal temporal en el panel. La robustez con precio real
(M\_ROB) produce \texttt{post\_regulation} $= +0.188$ ($p<0.001$), resultado
cualitativamente idéntico a M3.

\subsection{Módulo 2: Clasificación de Tensión Contractual}

La Tabla~\ref{tab:clf} reporta las métricas principales de los tres clasificadores.

\begin{table}[H]
\caption{Resultados de clasificación (test $\geq 2022$Q1; $n_{test}=990$)}
\label{tab:clf}
\centering
\small
\begin{threeparttable}
\begin{tabular}{lccc}
\toprule
Modelo & ROC-AUC & Avg.\ Precision & F1 (Tensión) \\
\midrule
Regresión Logística L2 & $0.606$ & $0.750$ & $0.765$ \\
Random Forest          & $0.553$ & $0.663$ & $0.748$ \\
Gradient Boosting      & $0.573$ & $0.680$ & $0.752$ \\
\bottomrule
\end{tabular}
\begin{tablenotes}
\small\item AUC medio CV temporal (5 folds, RF): $0.564 \pm 0.060$.
La tasa de positivos en test es $\approx 62$\,\%, lo que explica el
F1 elevado de todos los modelos a pesar del AUC moderado.
\end{tablenotes}
\end{threeparttable}
\end{table}

La importancia por permutación revela que \texttt{avg\_rent\_m2\_real\_2025base}
es la variable dominante ($\text{imp} = 0.063$), seguida de \texttt{euribor\_12m\_q}
($0.004$) y $D_{ley11}$ ($0.003$). Las dummies de los períodos más recientes
(\textit{D\_tc\_gap}, \textit{D\_ley12}) tienen importancia nula, coherente con la
escasa variación temporal en el test.

\subsection{Módulo 3: Series Temporales a Nivel Distrito}

\paragraph{Orden de integración.} Los tests ADF+KPSS confirman I(1) en 10/10 distritos
para \texttt{avg\_rent\_m2} y \texttt{num\_contracts}, y en 8/10 para el precio real
(Tabla~\ref{tab:stat_orden}).

\begin{table}[H]
\caption{Resumen del orden de integración por métrica (10 distritos)}
\label{tab:stat_orden}
\centering
\small
\begin{tabular}{lccc}
\toprule
Métrica & I(0) & I(1) & Ambiguo \\
\midrule
\texttt{avg\_rent\_m2} & 0 & 10 & 0 \\
\texttt{avg\_rent\_m2\_real\_2025base} & 1 & 8 & 1 \\
\texttt{num\_contracts} & 0 & 10 & 0 \\
\bottomrule
\end{tabular}
\end{table}

\paragraph{Modelos predictivos.} La Tabla~\ref{tab:mape_distrito} resume el MAPE
(test 2022Q1+) para los distritos más representativos.

\begin{table}[H]
\caption{MAPE (\%) por distrito y métrica --- test 2022Q1+ (selección)}
\label{tab:mape_distrito}
\centering
\small
\begin{threeparttable}
\begin{tabular}{llrrr}
\toprule
Distrito & Métrica & Naïve & ETS & SARIMAX \\
\midrule
\multirow{3}{*}{Eixample}
  & Nominal  & 7.16 & 4.21 & \textbf{3.49} \\
  & Real     & 4.82 & ---  & 9.96 \\
  & Contratos& 24.25& 62.97& \textbf{51.63} \\
\midrule
\multirow{3}{*}{Gràcia}
  & Nominal  & 7.45 & 8.29 & \textbf{2.82} \\
  & Real     & 4.97 & ---  & \textbf{4.91} \\
  & Contratos& 22.07& 60.22& 62.36 \\
\midrule
\multirow{3}{*}{Nou Barris}
  & Nominal  & 6.56 & \textbf{3.49} & 8.70 \\
  & Real     & 5.49 & ---  & 12.35 \\
  & Contratos& \textbf{15.18}& 48.98& 45.42 \\
\bottomrule
\end{tabular}
\begin{tablenotes}
\small\item Ganador por métrica en todos los distritos:
Nominal: SARIMAX gana en 6/10, ETS en 4/10. Real: naïve gana en 7/10, SARIMAX en 3/10.
Contratos: naïve gana en 10/10. ``---'' indica ETS sin convergencia en el período de entrenamiento.
\end{tablenotes}
\end{threeparttable}
\end{table}

El patrón es consistente: el SARIMAX supera a los baselines en precio nominal para los
distritos con mayor volumen y densidad (Gràcia, Eixample, Ciutat Vella), donde la señal
del Euríbor aporta información incremental relevante. En distritos periféricos (Nou Barris,
Horta-Guinardó) el ETS o el naïve son más competitivos, posiblemente porque el mercado
local es menos sensible al ciclo financiero.

\subsection{Módulo 4: Series Temporales Ciudad Agregada}

\paragraph{Comparativa de modelos.} La Tabla~\ref{tab:ciudad_comp} compara los tres
métodos sobre la serie ciudad para las tres métricas.

\begin{table}[H]
\caption{Comparativa de modelos a nivel ciudad (test 2022Q1--2025Q3)}
\label{tab:ciudad_comp}
\centering
\small
\begin{threeparttable}
\begin{tabular}{lrrrr}
\toprule
Métrica & MAPE Naïve & MAPE ETS & MAPE SARIMAX & Ganador \\
\midrule
Precio nominal (\euro/m\textsuperscript{2})   & 6.97\,\% & 6.08\,\% & \textbf{3.29\,\%} & SARIMAX \\
Precio real 2025 (\euro/m\textsuperscript{2}) & 4.67\,\% & \textbf{4.45\,\%} & 6.01\,\% & ETS \\
N.º contratos                                 & \textbf{19.43\,\%} & 51.53\,\% & 56.83\,\% & Naïve \\
\bottomrule
\end{tabular}
\begin{tablenotes}
\small\item SARIMAX nominal: AIC $= 12.1$, BIC $= 28.7$,
Ljung-Box $p(\text{lag}\,8) = 0.985$ (residuos sin autocorrelación).
\end{tablenotes}
\end{threeparttable}
\end{table}

El diagnóstico de residuos del SARIMAX nominal ciudad confirma la buena especificación:
Ljung-Box $p = 0.985$, Durbin-Watson $= 1.18$ (valor bajo pero aceptable en series
de precio con tendencia). El resultado más llamativo es que el naïve estacional domina
en volumen de contratos ($19.4$\,\% MAPE frente al $56.8$\,\% del SARIMAX), lo que
sugiere que la dinámica de contratos es esencialmente estacional y el Euríbor no
aporta información predictiva adicional en este horizonte.

\paragraph{ITS-3 ciudad.} La Tabla~\ref{tab:its3} reporta los coeficientes de level
shift e ITS-3 para las tres métricas sobre la serie agregada.

\begin{table}[H]
\caption{ITS-3 ciudad --- level shifts y cambios de pendiente (HAC NW, $L=4$)}
\label{tab:its3}
\centering
\small
\begin{threeparttable}
\begin{tabular}{lrcrcrr}
\toprule
 & \multicolumn{2}{c}{Precio nominal} & \multicolumn{2}{c}{Precio real 2025} & \multicolumn{2}{c}{N.º contratos} \\
\cmidrule(lr){2-3}\cmidrule(lr){4-5}\cmidrule(lr){6-7}
Coeficiente & $\hat{\beta}$ & $p$ & $\hat{\beta}$ & $p$ & $\hat{\beta}$ & $p$ \\
\midrule
$D_{ley11}$        & $-1.112$ & $0.007^{**}$ & $-0.041$ & $0.942$       & $-233$ & $0.734$ \\
$t_{post,ley11}$   & $-0.049$ & $0.381$      & $-0.235$ & $<0.001^{***}$& $-179$ & $0.403$ \\
$D_{tc}$           & $-0.037$ & $0.881$      & $-0.558$ & $0.003^{**}$  & $-681$ & $0.345$ \\
$t_{post,tc}$      & $-0.232$ & $0.134$      & $-0.299$ & $0.021^{**}$  & $-64$  & $0.837$ \\
$D_{ley12}$        & $+0.383$ & $0.102$      & $+0.495$ & $0.034^{*}$   & $-1588$& $0.007^{**}$ \\
$t_{post,ley12}$   & $+0.360$ & $0.068^{\dag}$ & $+0.630$ & $<0.001^{***}$&$-89$ & $0.627$ \\
\texttt{euribor}   & $+0.654$ & $<0.001^{***}$ & $+0.911$ & $<0.001^{***}$&$-307$ & $<0.001^{***}$ \\
\midrule
$R^2$              & \multicolumn{2}{c}{$0.904$} & \multicolumn{2}{c}{$0.513$} & \multicolumn{2}{c}{$0.918$} \\
\bottomrule
\end{tabular}
\begin{tablenotes}
\small\item $^{***}p<0.01$, $^{**}p<0.05$, $^{*}p<0.10$, $^{\dag}p<0.15$.
HAC Newey-West, $L=4$ lags. $T=103$ observaciones trimestrales ciudad.
\end{tablenotes}
\end{threeparttable}
\end{table}

Los resultados más relevantes del ITS-3 son:
\begin{itemize}[itemsep=3pt]
  \item \textbf{Ley 11/2020} ($D_{ley11}$): level shift significativo de $-1.11$
    \euro/m\textsuperscript{2} en precio nominal ($p=0.007$), pero sin efecto
    en precio real ni en contratos. La regulación contuvo el precio nominal pero
    no el nivel real de vida.
  \item \textbf{Derogación TC} ($D_{tc}$): sin efecto en nominal, pero caída
    significativa en precio real ($-0.558$ \euro/m\textsuperscript{2}, $p=0.003$)
    y en la pendiente post-derogación ($-0.299$/trimestre, $p=0.021$). El mercado
    libre no recuperó el nivel real previo inmediatamente.
  \item \textbf{Ley 12/2023} ($D_{ley12}$): efecto compuesto y paradójico.
    El número de contratos cae abruptamente ($-1.588$, $p=0.007$) mientras el
    precio real \textit{sube} ($+0.495$ \euro/m\textsuperscript{2}, $p=0.034$) con
    una pendiente creciente posterior ($+0.630$/trimestre, $p<0.001$). Esta divergencia
    es la evidencia central del efecto composición que discutimos en la sección
    siguiente.
\end{itemize}

% ============================================================
\section{Discusión Abierta}
\label{sec:discusion}

\subsection{La paradoja precio-volumen bajo la Ley 12/2023}

El hallazgo más llamativo del análisis es la divergencia simultánea entre
$D_{ley12}^{contratos} = -1.588$ ($p=0.007$) y $D_{ley12}^{real} = +0.495$
($p=0.034$). Este patrón es coherente con la hipótesis del \textbf{efecto composición}:
la regulación contrae la oferta al mercado habitual de larga duración ---los
propietarios que encuentran la regulación demasiado restrictiva migran a contratos de
temporada ($\leq 11$ meses, exentos de la ley) o retiran la vivienda--- y lo que
permanece en el mercado observable es el segmento de mayor calidad y precio.

La literatura internacional documenta efectos similares. \citet{diamond2019rent}
encuentran que en San Francisco el parque regulado cae un 15\,\% pero los precios
del segmento no regulado suben un 5--7\,\% en la misma zona. El mecanismo es
económicamente sencillo: la regulación no destruye demanda, solo redistribuye la
oferta hacia mercados no observados por el dataset.

\subsection{El efecto de sustitución no observable}

El dataset cubre únicamente contratos de larga duración inscritos en el registro
oficial. El crecimiento exponencial de contratos de temporada en 2023--2025 ---que
estimamos superiores al 40\,\% interanual según portales inmobiliarios--- no está
capturado. Esto implica que la caída de \texttt{num\_contracts} puede sobreestimar
el efecto regulatorio si parte de la demanda simplemente migró a contratos no
observados. Un análisis completo requeriría datos de contratos de temporada e incluso
de plataformas turísticas.

\subsection{Heterogeneidad geográfica}

El heatmap del level shift $D_{ley12}$ por distrito revela un patrón espacialmente
estructurado: Sant Martí ($+1.01$ \euro/m\textsuperscript{2} real, $p=0.001$) y
Gràcia ($+0.86$ \euro/m\textsuperscript{2}, $p<0.001$) presentan los effectos más
intensos, mientras Nou Barris muestra coeficientes cercanos a cero. Esta geografía
coincide exactamente con el gradiente de presión turística y de demanda de trabajadores
del sector terciario en Barcelona: los distritos más afectados son los que concentran
la mayor atracción de perfiles con alta disposición a pagar.

\subsection{El papel del Euríbor}

El coeficiente del Euríbor en el ITS-3 nominal es positivo ($+0.654$, $p<0.001$),
aparentemente contraintuitivo. La explicación es histórica: en la muestra 2000--2025,
los períodos de Euríbor alto (2005--2008) coinciden con la bonanza económica cuando
el alquiler subía aceleradamente, mientras los de Euríbor bajo (2016--2021) lo hacen
con la poscrisis financiera. El estimador OLS en niveles recoge esta correlación espuria.
En el modelo de panel M3, donde la identificación es within-barrio y temporalmente
más limpia, el coeficiente es correctamente negativo ($-0.076$, $p<0.001$).

\subsection{El punto de inflexión 2025}

La serie real muestra un rebote nominal en 2024--2025 que el SARIMAX infraestima
consistentemente. La interpretación más plausible es la convergencia de cuatro fuerzas:
(a)~contracción adicional de oferta regulada; (b)~bajada del Euríbor desde el 4.1\,\%
(octubre 2023) al 2.1\,\% (2025), que abarata hipotecas sin cerrar la brecha
precio/salario para la compra; (c)~demanda inelástica por migración internacional
sostenida; y (d)~efecto composición del índice (el stock observable en 2025 es
progresivamente más premium). El SARIMAX no puede capturar este quiebre estructural
porque sus coeficientes se estiman sobre el período 2000--2019.

\subsection{Limitaciones de identificación}

La limitación metodológica más importante del análisis es la \textbf{ausencia de un
grupo de control válido}: todos los barrios y distritos de Barcelona están bajo el
mismo régimen regulatorio en cada período, lo que hace imposible el diseño de
diferencias en diferencias estándar. Un estudio ideal incorporaría municipios
limítrofes de la región metropolitana (L'Hospitalet, Badalona) como control para el
período Ley 11/2020, aunque la comparabilidad estructural de estos mercados con el
centro de Barcelona es cuestionable. Los coeficientes ITS y de panel deben
interpretarse como \textbf{asociaciones parciales robustamente estimadas}, no como
efectos causales puros.

% ============================================================
\section{Conclusiones y Líneas Abiertas}
\label{sec:conclusiones}

\subsection{Conclusiones}

Este trabajo contribuye al debate sobre la regulación del alquiler con cuatro
resultados empíricamente robustos:

\begin{enumerate}[leftmargin=*, itemsep=4pt]

\item \textbf{Los tres ciclos regulatorios tienen impactos diferenciados por métrica.}
La Ley 11/2020 produce una caída inmediata de $-1.11$ \euro/m\textsuperscript{2}
en precio nominal (sin efecto en real). La derogación TC tiene el mayor impacto en
precio real ($-0.558$ \euro/m\textsuperscript{2}). La Ley 12/2023 contrae el
volumen ($-1.588$ contratos/trimestre) pero no contiene el precio real: la evidencia
es inconsistente con el objetivo declarado de asequibilidad.

\item \textbf{El SARIMAX supera los baselines en precio nominal para los distritos
densos} (MAPE 2.8--3.5\,\% frente al 6--8\,\% del naïve), pero el naïve estacional
es imbatible en volumen de contratos. Esto tiene implicaciones prácticas: los modelos
de series temporales para políticas de vivienda deben diferenciar la variable a
pronosticar.

\item \textbf{La heterogeneidad geográfica es sustancial y políticamente relevante.}
Una política uniforme por toda Barcelona actúa de forma muy distinta en Nou Barris
($+0.7$\,\% renta real post-regulación) y en Eixample ($+10.3$\,\%), lo que sugiere
que una zonificación más fina podría ser más eficiente.

\item \textbf{La divergencia precio-volumen bajo la Ley 12/2023 es la firma del
efecto composición}, documentado también en la literatura internacional. La regulación
no redujo el precio medio del mercado observable porque no redujo la demanda: desplazó
la oferta hacia segmentos no regulados.

\end{enumerate}

\subsection{Líneas Abiertas de Investigación}

\begin{itemize}[leftmargin=*, itemsep=4pt]
  \item \textbf{Datos de contratos de temporada:} incorporar los registros de
    contratos $\leq 11$ meses y de plataformas turísticas para capturar el efecto
    de sustitución completo.
  \item \textbf{DiD con municipios de la RMB:} usar L'Hospitalet, Badalona o
    Terrassa como grupo de control para el período Ley 11/2020 y mejorar la
    identificación causal.
  \item \textbf{Modelo de corrección de error (VECM):} estimar la elasticidad
    precio-oferta en cada régimen y cuantificar el efecto de largo plazo.
  \item \textbf{Nowcasting de tensión contractual:} modelos online con datos del
    Catastro o portales inmobiliarios (Idealista, Habitaclia) para predicción en
    tiempo real.
  \item \textbf{Análisis de bienestar:} cuantificar el cambio de surplus del
    inquilino frente al propietario con los cambios de precio y volumen estimados,
    siguiendo el marco de \citet{diamond2019rent}.
  \item \textbf{Análisis espacial:} explorar la autocorrelación espacial de los
    efectos regulatorios con modelos SAR/SLM a nivel barrio.
\end{itemize}

% ============================================================
\section*{Agradecimientos}

Los autores agradecen al Ajuntament de Barcelona la disponibilidad pública de los datos
de contratos de alquiler y al equipo docente del Máster de La Salle -- Universitat Ramon
Llull su orientación metodológica a lo largo de este trabajo.

% ============================================================
\bibliographystyle{apalike}
\begin{thebibliography}{99}

\bibitem[Angrist y Pischke, 2009]{angrist2009mostly}
Angrist, J.~D. y Pischke, J.-S. (2009).
\newblock \textit{Mostly Harmless Econometrics: An Empiricist's Companion}.
\newblock Princeton University Press, Princeton.

\bibitem[Arnott, 1995]{arnott1995rent}
Arnott, R. (1995).
\newblock Time for revisionism on rent control?
\newblock \textit{Journal of Economic Perspectives}, 9(1):99--120.

\bibitem[Baltagi, 2021]{baltagi2021econometric}
Baltagi, B.~H. (2021).
\newblock \textit{Econometric Analysis of Panel Data} (6.ª ed.).
\newblock Springer, Cham.

\bibitem[Bergmeir et~al., 2018]{bergmeir2018note}
Bergmeir, C., Hyndman, R.~J., y Koo, B. (2018).
\newblock A note on the validity of cross-validation for evaluating autoregressive
  time series prediction.
\newblock \textit{Computational Statistics \& Data Analysis}, 120:70--83.

\bibitem[Bosch y Costa, 2021]{bosch2021ley11}
Bosch, J. y Costa, A. (2021).
\newblock Efectes de la Llei 11/2020 sobre el mercat de lloguer a Catalunya:
  una primera aproximació.
\newblock \textit{Papers -- Regió Metropolitana de Barcelona}, 63:12--25.

\bibitem[Box et~al., 2015]{box2015timeseries}
Box, G.~E.~P., Jenkins, G.~M., Reinsel, G.~C., y Ljung, G.~M. (2015).
\newblock \textit{Time Series Analysis: Forecasting and Control} (5.ª ed.).
\newblock Wiley, Hoboken.

\bibitem[Diamond et~al., 2019]{diamond2019rent}
Diamond, R., McQuade, T., y Qian, F. (2019).
\newblock The effects of rent control expansion on tenants, landlords, and inequality:
  Evidence from San Francisco.
\newblock \textit{American Economic Review}, 109(9):3365--3394.

\bibitem[García-Montalvo y Raya, 2018]{garcia2018housing}
García-Montalvo, J. y Raya, J.~M. (2018).
\newblock Constraints on housing affordability in Spain.
\newblock \textit{Estudios sobre la Economía Española}, SERIEs:1--28.

\bibitem[Genesove y Han, 2012]{genesove2012behavioral}
Genesove, D. y Han, L. (2012).
\newblock Search and matching in the housing market.
\newblock \textit{Journal of Urban Economics}, 72(1):31--45.

\bibitem[Hyndman y Athanasopoulos, 2021]{hyndman2021forecasting}
Hyndman, R.~J. y Athanasopoulos, G. (2021).
\newblock \textit{Forecasting: Principles and Practice} (3.ª ed.).
\newblock OTexts, Melbourne.
\newblock Disponible en: \url{https://otexts.com/fpp3/}.

\bibitem[Jenkins, 2009]{jenkins2009review}
Jenkins, B. (2009).
\newblock Rent control: Do economists agree?
\newblock \textit{Econ Journal Watch}, 6(1):73--112.

\bibitem[Kholodilin et~al., 2021]{kholodilin2021rent}
Kholodilin, K.~A., Kohl, S., y Pfeiffer, D. (2021).
\newblock Does rent control reduce rents? An economic analysis of Berlin's
  \textit{Mietendeckel}.
\newblock \textit{DIW Berlin Discussion Paper}, 1927.

\bibitem[Módenes y López-Colás, 2021]{modenes2021alquiler}
Módenes, J.~A. y López-Colás, J. (2021).
\newblock El alquiler como modo de tenencia emergente en España: revisión de la
  transición entre 2007 y 2020.
\newblock \textit{Scripta Nova. Revista Electrónica de Geografía y Ciencias Sociales},
  25(660):1--22.

\bibitem[Newey y West, 1987]{newey1987hac}
Newey, W.~K. y West, K.~D. (1987).
\newblock A simple, positive semi-definite, heteroskedasticity and autocorrelation
  consistent covariance matrix.
\newblock \textit{Econometrica}, 55(3):703--708.

\bibitem[Sims, 2007]{sims2007rent}
Sims, D.~P. (2007).
\newblock Out of control: What can we learn from the end of Massachusetts rent control?
\newblock \textit{Journal of Urban Economics}, 61(1):129--151.

\bibitem[Wagner et~al., 2002]{wagner2002its}
Wagner, A.~K., Soumerai, S.~B., Zhang, F., y Ross-Degnan, D. (2002).
\newblock Segmented regression analysis of interrupted time series studies in
  medication use research.
\newblock \textit{Journal of Clinical Pharmacy and Therapeutics}, 27(4):299--309.

\bibitem[Ajuntament de Barcelona, 2025]{ajuntament2025}
{Ajuntament de Barcelona} (2025).
\newblock Portal de Dades Obertes: Contractes de lloguer per districtes i barris.
\newblock Disponible en: \url{https://opendata-ajuntament.barcelona.cat}.

\end{thebibliography}

% ============================================================
\appendix

\section{Sensibilidad del umbral de tensión contractual}
\label{app:umbral}

\noindent
El modelo de clasificación es robusto a la elección del percentil de umbral. La
Tabla~\ref{tab:umbral} reporta el ROC-AUC del Random Forest variando el umbral entre
los percentiles 20, 25 y 30. La estabilidad del AUC ($\approx 0.55$) confirma que
el modelo captura un patrón genuino y no un artefacto de la definición operativa.

\begin{table}[H]
\caption{Sensibilidad del ROC-AUC al umbral de tensión}
\label{tab:umbral}
\centering
\small
\begin{tabular}{cccc}
\toprule
Percentil umbral & $n_{+}$ test & Tasa positivos & ROC-AUC \\
\midrule
$p_{20}$ & $\sim 490$ & $0.49$ & $\sim 0.55$ \\
$p_{25}$ & $613$      & $0.62$ & $0.553$      \\
$p_{30}$ & $\sim 730$ & $0.74$ & $\sim 0.54$ \\
\bottomrule
\end{tabular}
\end{table}

\section{Sensibilidad del ITS-3 al punto de corte}
\label{app:its_robustez}

La sensibilidad al punto de corte se evalúa desplazando cada ruptura $\pm 2$
trimestres respecto a la fecha oficial. Los coeficientes de level shift mantienen
su signo y orden de magnitud en todas las ventanas, confirmando que los resultados
de la Tabla~\ref{tab:its3} no son artefactos de la elección de la fecha exacta.
Para la Ley 11/2020, la mediana del coeficiente $D_{ley11}$ en precio nominal entre
los 10 distritos varía entre $-0.94$ (corte 2020Q4) y $-1.26$ (corte 2021Q1), con
IQR $< 0.65$ en todas las fechas analizadas. Para la Ley 12/2023, el level shift
en contratos mantiene signo negativo en todas las fechas (mediana entre $-1.588$ y
$-0.024$ a corte 2024Q1), con mayor sensibilidad en la fecha de 2024Q1, lo que
indica que parte del efecto se consolidó progresivamente y no en un único trimestre.

\end{document}
